\section{Abréviations du carnet}

Hippolyte Combes écrit au jour le jour, souvent en marche ou à cheval, et abrège sans cesse : l'édition y a relevé quelque 950 formes abrégées, pour plus de 3\,000 occurrences. L'apparat ne développe que les formes rares ou ambiguës ; les plus fréquentes sont réunies ici. Son système repose sur trois procédés, ceux de l'usage manuscrit courant : la contraction, qui garde le début et la fin du mot et place volontiers la finale en lettres suscrites, presque jamais suivie d'un point ({\LJpolice M\LJsup{me}, appart\LJsup{t}, tempér\LJsup{re}}), et dont relève la série très productive des adverbes en \textit{-t} ({\LJpolice ext\LJsup{t}, mervt, absolt}) ; la suspension, qui coupe le mot et le fait souvent suivre d'un point ({\LJpolice Sam., égl., Cap.}) ; enfin des signes et des chiffres employés comme mots ({\LJpolice \&}, \textit{et} ; {\LJpolice +}, \textit{plus} ; {\LJpolice =}, \textit{même} ; {\LJpolice 1}, \textit{un} ; {\LJpolice s/}, \textit{sur}). À l'armée britannique, il emprunte en outre, tels quels, ses sigles et ses abréviations d'unités ({\LJpolice Coy, A.S.C., R.T.O.}). La table donne la forme diplomatique, le développement en italique et le nombre d'occurrences relevées pour l'ensemble de la ligne ; les variantes d'une même abréviation sont séparées par une virgule, les mots distincts par un point-virgule, et «~?~» signale un développement non assuré.

\begingroup
\footnotesize
\setlength{\tabcolsep}{3pt}
\renewcommand{\arraystretch}{1.06}
\setlength{\LTpre}{6pt}
\setlength{\LTpost}{0pt}
\begin{longtable}{@{}>{\LJpolice\raggedright\arraybackslash}p{51.5mm}>{\itshape\raggedright\arraybackslash}p{76.5mm}>{\raggedleft\arraybackslash}p{7mm}@{}}
\toprule
\multicolumn{1}{@{}l}{\normalfont\bfseries Forme} & \multicolumn{1}{l}{\normalfont\bfseries Développement} & \multicolumn{1}{r@{}}{\normalfont\bfseries Occ.}\\
\midrule
\endfirsthead
\multicolumn{3}{@{}l}{\normalfont\itshape Abréviations du carnet (suite)}\\
\toprule
\multicolumn{1}{@{}l}{\normalfont\bfseries Forme} & \multicolumn{1}{l}{\normalfont\bfseries Développement} & \multicolumn{1}{r@{}}{\normalfont\bfseries Occ.}\\
\midrule
\endhead
\bottomrule
\endfoot
\bottomrule
\endlastfoot
\multicolumn{3}{@{}l}{\normalfont\bfseries Jours de la semaine}\\*
\addlinespace[1pt]
L, L. & lundi & 11\\
Ma, ma, Ma., Mar, Ma\LJsup{rdi} & mardi & 17\\
Me, me, Me., Me –, mer, mer., mer-, Mer., Mer-, Merc, merc, Merc., merc., merc-, merc– & mercredi & 36\\
J, J. & jeudi & 12\\
Ve, Ve., Ven, Ven., Ven–, Vendr, Vendr.~; 1\LJsup{er} V. & vendredi & 32\\
S., Sam, Sam. & samedi \textup{(S. aussi Saint)} & 33\\
D, D., Dim, Dim. & dimanche & 36\\
\addlinespace[3pt]
\multicolumn{3}{@{}l}{\normalfont\bfseries Mois et millésimes}\\*
\addlinespace[1pt]
Sept, sept, Sept\LJsup{re}, Sept\LJsup{r}, Sept\LJsup{bre}, Septb, Septbr, Septbre, Sepbr & septembre & 26\\
Oct, oct, oct., Oct\LJsup{b}, Oct\LJsup{bre}, oct\LJsup{bre}, octbr, Octbre & octobre & 28\\
Nov, Nov., Nov –, Nov\LJsup{b}, Nov\LJsup{bre}, Novbre & novembre & 10\\
Dec, Dec\LJsup{b}, Decb, Dec\LJsup{bre}, Déc.\LJsup{b}, Déc\LJsup{bre}, Décbre & décembre & 11\\
Janv, Janv. & janvier & 6\\
Fev., Fév, Fév-, Fevr, Fevr. & février & 7\\
/16, 16~; /17, 17, '17 & 1916~; 1917 & 78\\
\addlinespace[3pt]
\multicolumn{3}{@{}l}{\normalfont\bfseries Mots grammaticaux et mots courants}\\*
\addlinespace[1pt]
\& & et & 386\\
+~; de + en +~; +tot~; La + & plus~; de plus en plus~; plutôt~; la plupart & 31\\
=, moi-=, de = & même & 3\\
1, 1 peu, d'1, l'1, à 1~; les 1 & un, une~; les uns & 105\\
ds, Ds & dans & 131\\
ns, ns., Ns & nous & 20\\
s/ & sur & 54\\
pr, pr., p. & pour & 17\\
av & avec & 30\\
chz, Chz, ch & chez & 29\\
cme, Cme, co—e & comme & 14\\
tt, Tt~; tte, Tte~; ttes~; ts, Ts & tout~; toute~; toutes~; tous & 85\\
qq, Qq~; ql. mnts & quelque, quelques~; quelques moments \textup{(ou minutes ?)} & 37\\
bp & beaucoup & 12\\
tr, tre, Tre, trè, Trè & très & 9\\
asz, assz & assez & 14\\
plusrs, Plusrs, plusis, plus., pl. & plusieurs & 10\\
presq~; jusq~; pt~; e. & presque~; jusque~; peut~; encore & 4\\
b., b & belle, belles~; bonne \textup{(de b heure)} & 9\\
c.~; l., sur l.~; j. & certain(s)~; laquelle~; jeune & 9\\
d & de \textup{(d l'office, d St Hippte)} & 2\\
gd, Gd, g\LJsup{d}, gr, grd~; gde, g\LJsup{de}~; gds, g\LJsup{ds}, gdes & grand~; grande~; grands, grandes \textup{(aussi gd-mère, g\LJsup{d}-route, gd-place)} & 53\\
tjs, Tjs~; Tjours, tjours, toujrs & toujours & 31\\
lgs, lg, lgt, lgt., longt & longtemps & 30\\
aujd'hui, Aujd'hui, aujdhui, auj'dhui, Auj'hui, Auj'd'hui, Auj-dhui & aujourd'hui & 12\\
j-à, J-à, j'à, jà, j à, j.à, j. à, jusq à~; j-au, jus, j~; j-aux~; j-après & jusqu'à~; jusqu'au~; jusqu'aux~; jusqu'après & 48\\
à g., A g., g.~; à d., dr. & à gauche~; à droite & 11\\
côt~; Aprè-midi & côté~; après-midi & 3\\
etc, etc., etc. etc. & et cetera & 33\\
", id. & idem & 3\\
Cf. & confer & 1\\
Y & il y a & 1\\
x \textnormal{(de x)} & \textup{nombre indéterminé :} de nombreux & 1\\
\addlinespace[3pt]
\multicolumn{3}{@{}l}{\normalfont\bfseries Adverbes en -ment}\\*
\addlinespace[1pt]
ext, ext\LJsup{t}, extt, extr., extr\LJsup{t}, Extremt, extrêm\LJsup{t} & extrêmement & 24\\
merv\LJsup{t}, merv.\LJsup{t}, mervt, merveillt, merveill\LJsup{t} & merveilleusement & 12\\
affr\LJsup{t}, affrt, affreut, affreust & affreusement & 4\\
absol\LJsup{t}, absolt, abst, abt & absolument & 6\\
admir\LJsup{t}., Admirab\LJsup{t}~; aimablt~; visib\LJsup{t} & admirablement~; aimablement~; visiblement & 4\\
entt, entiert, entier\LJsup{t} & entièrement & 8\\
fort\LJsup{t}, fortt~; rapid\LJsup{t}, rapidt & fortement~; rapidement & 9\\
definit, definit\LJsup{t}, définitt & définitivement & 3\\
comp\LJsup{t}, compt, complet\LJsup{t}, complèt\LJsup{t} & complètement & 4\\
délic.\LJsup{t}, delict, delicieus\LJsup{t} & délicieusement & 3\\
principt~; magnif\LJsup{t}~; considt~; lamentablt & principalement~; magnifiquement~; considérablement~; lamentablement & 5\\
exceptt~; excessv\LJsup{t}~; infinit~; notam~; abond\LJsup{t} & exceptionnellement~; excessivement~; infiniment~; notamment~; abondamment & 5\\
inutilt~; profondt~; plaintivt~; parallelt~; seulemt~; lentemt & inutilement~; profondément~; plaintivement~; parallèlement~; seulement~; lentement & 6\\
\addlinespace[3pt]
\multicolumn{3}{@{}l}{\normalfont\bfseries Substantifs et adjectifs contractés ou suspendus}\\*
\addlinespace[1pt]
appart\LJsup{t}, appartt~; compart\LJsup{t}, compartt, compartimt, Compartimt & appartement~; compartiment & 15\\
bombard\LJsup{t}, bombardt, bombardemt. & bombardement & 5\\
changem, changemt~; comm\LJsup{t}, Comm\LJsup{t}, commt~; mouv\LJsup{t}, mouvt & changement~; commencement~; mouvement & 7\\
derangt~; arrangt~; attachemt~; gémisst~; sentimt~; epanchts~; enchevêtrt & dérangement~; arrangement~; attachement~; gémissement~; sentiment~; épanchements~; enchevêtrement & 7\\
renseig\LJsup{ts}, renseigts~; enregt, Enregt~; cantonnem\LJsup{t}, cantonnts & renseignements~; enregistrement~; cantonnement(s) & 7\\
dep\LJsup{t}, départt~; rég\LJsup{t} & département~; régiment & 3\\
prop\LJsup{re}, propr\LJsup{e}, propre, propri, proprie, Proprie~; proprios & propriétaire~; propriétaires & 13\\
mil\LJsup{re}, militre~; ord\LJsup{re}, ordin\LJsup{re}~; annivre~; comm\LJsup{re} & militaire~; ordinaire~; anniversaire~; commissionnaire & 8\\
auxil\LJsup{re}~; celib\LJsup{re}~; sous-pref\LJsup{re}~; permiss\LJsup{re}, permissre & auxiliaire~; célibataire~; sous-préfecture~; permissionnaires & 5\\
tempér\LJsup{re}, Tempér\LJsup{re}, Tempér\LJsup{e}, tempérre, Tempérre~; chal\LJsup{r} & température~; chaleur & 11\\
fr, fr., fr\LJsup{e}, Fr.\LJsup{se}, fr\LJsup{se}, F\LJsup{se}, f\LJsup{se}, frse, frsé, fse, Fse, F\LJsup{r}, F\LJsup{r}.~; frs, frç~; anglse & française~; français~; anglaise & 30\\
cath\LJsup{le}, cathle~; horiz\LJsup{le}~; artif\LJsup{lles}~; abomibles & cathédrale~; horizontale~; artificielles~; abominables & 8\\
merv\LJsup{x}, Merv\LJsup{x}, Mervx~; merv\LJsup{se}~; affr. & merveilleux~; merveilleuse~; affreux & 7\\
bienh\LJsup{x}~; nouvx~; chevx~; jx~; Gatx & bienheureux~; nouveaux~; chevaux~; journaux~; gâteaux & 5\\
exc\LJsup{te}, Exc\LJsup{te}~; imptes~; verdoyte~; {\NoAutoSpacing intermitt:}~; extre & excellente~; importantes~; verdoyante~; intermittente(s)~; extrême & 6\\
symphe, sympathiq\LJsup{ue}~; pittoresq~; elastiq~; nuag~; eff\LJsup{l} & sympathique~; pittoresques~; élastiques~; nuageux~; effilé & 6\\
Corresp\LJsup{ce}, correspce~; confer\LJsup{ce}~; intellige~; intermitt\LJsup{ce}~; reconnaissce & correspondance~; conférence~; intelligence~; intermittence~; reconnaissance & 9\\
convers\LJsup{n}, Conver\LJsup{it}.~; sit\LJsup{on}~; autorison~; confirmaton~; install° & conversation~; situation~; autorisation~; confirmation~; installation & 7\\
réparatn, repar~; Assomptn~; progro— & réparation~; Assomption~; programme & 4\\
100ne~; douz. & centaine~; douzaines & 2\\
cie, Cie & compagnie & 10\\
divi, Div\LJsup{n}, divon~; Div., div. & division~; diverses & 7\\
égl., egl~; Hoes~; fem. & église~; hommes~; femmes & 4\\
cl., Cl. & classe & 10\\
F\LJsup{me}~; chps~; chamb~; Bd, B\LJsup{d}~; carref & ferme~; champs~; chambre~; boulevard~; carrefour & 8\\
ét.~; rez-de-ch.~; bag.~; voy.~; voy\LJsup{rs}~; hab\LJsup{ts}~; souvrs & étage~; rez-de-chaussée~; bagages~; voyages~; voyageurs~; habitants~; souvenirs & 7\\
choc\LJsup{t}~; diam~; Trams~; Tre & chocolat~; diamant~; tramways~; teinture \textup{(Tre d'iode)} & 7\\
dejeun, dejeun., déjeun, Dejeun~; Lev. & déjeuner, déjeuné~; levé & 8\\
renseig~; surchauff \textnormal{(en fin de ligne)} & renseigner~; surchauffée & 2\\
Mag, Mag.~; Chge & magasins~; change & 3\\
Ptit~; l'extré~; long\LJsup{ues}~; à l – & petit~; extrémité~; longues~; à lettres & 4\\
Géog., geog.~; Exerc.~; prof., prof & géographie, géographiques~; exercices~; professeur & 5\\
M-Age~; s. \textnormal{(XIII\LJsup{e} s.)} & Moyen Âge~; siècle & 2\\
Norm\LJsup{le} Sup\LJsup{re}~; Poly\LJsup{te} & \textup{(École)} normale supérieure~; \textup{(École)} polytechnique & 2\\
\addlinespace[3pt]
\multicolumn{3}{@{}l}{\normalfont\bfseries Titres, grades, fonctions}\\*
\addlinespace[1pt]
M., M –~; M\LJsup{me}, Mme~; Mlle & Monsieur~; Madame~; Mademoiselle & 54\\
Ve – — & veuve \textup{(?)} & 1\\
M\LJsup{e}~; Dr.~; P.~; P.P. & Maître \textup{(notaire)}~; docteur~; Père~; Pères & 4\\
St, St., S., S\LJsup{t}~; Ste, S\LJsup{te} & Saint~; Sainte \textup{(aussi dans les noms de lieux)} & 56\\
Lt, Lt., L\LJsup{t}, l\LJsup{t}, Lieut, Lieu\LJsup{t}, lieut\LJsup{t}, Lieut\LJsup{t} & lieutenant & 36\\
Sec\LJsup{d} Lt & \textit{Second Lieutenant} \textup{(sous-lieutenant)} & 1\\
cap, Cap, cap., Cap., cap\LJsup{e}, cap\LJsup{ne}, Cap\LJsup{ne} & capitaine & 40\\
Capt., Cap\LJsup{t} & \textit{Captain} & 2\\
Col. & colonel & 1\\
gal, gen\LJsup{l}, genl & général & 3\\
marl, mal \textnormal{(devant} logis\textnormal{)} & maréchal \textup{(des logis)} & 3\\
int\LJsup{te}, intte, interp\LJsup{te}, interpte, intpte & interprète & 7\\
off, off., offs, offs., offrs, off\LJsup{r} & officier, officiers & 28\\
\addlinespace[3pt]
\multicolumn{3}{@{}l}{\normalfont\bfseries Heures, mesures, ordinaux}\\*
\addlinespace[1pt]
9\LJsup{h}, 7\LJsup{H}, 9 H, 6H, 2 h, 2 h., h. & heure, heures & 219\\
m, m.~; s, s. \textnormal{(après l'heure)} & du matin~; du soir & 73\\
½, 1/2~; au ½~; ½ jour & demi, et demie~; au milieu~; demi-jour, midi & 80\\
¼, 1/4 & quart & 4\\
m, m. \textnormal{(après un nombre)} & minutes & 11\\
m. \textnormal{(2m. 40, 50 m.)} & mètres & 3\\
K\LJsup{ils}, kil, kilom.~; kos, kos.~; cm~; l. & kilomètres~; kilos~; centimètres~; litres & 8\\
1\LJsup{er}, 1er, 1ers~; 1\LJsup{re}, 1\LJsup{ère}, 1ère & premier, premiers~; première & 20\\
2\LJsup{e}, 2e, 4\LJsup{e}, 54\LJsup{e}, 18\LJsup{ie}~; XIII\LJsup{e}, XIIIe, V\LJsup{e}, 5\LJsup{me} & deuxième, quatrième…~; treizième, cinquième \textup{(corps d'armée, armée)} & 22\\
no, n\LJsup{os} & numéro, numéros & 9\\
ct, ct. & courant \textup{(du mois)} & 4\\
\addlinespace[3pt]
\multicolumn{3}{@{}l}{\normalfont\bfseries Monnaie et comptes}\\*
\addlinespace[1pt]
1\LJsup{f}, 29\LJsup{f}, 2 f, 3f, 1fc.~; 1 fr, 84 fr., fr. & franc, francs & 98\\
60c, 60 c, 20c., 25 c, 75\LJsup{c} & centimes & 76\\
90., 60.. & centimes \textup{(unité non écrite)} & 2\\
C\LJsup{tes} p\LJsup{les}, c\LJsup{tes} p\LJsup{les}, Ctes ples, ctes ples, Cartes ples, Carte ple, C. p\LJsup{les}, C. P~; ples, p\LJsup{les}~; Cte & cartes postales~; postales~; carte & 54\\
\addlinespace[3pt]
\multicolumn{3}{@{}l}{\normalfont\bfseries Noms de personnes et de lieux}\\*
\addlinespace[1pt]
Marg\LJsup{te}, Margte, Marg., M\LJsup{te}, Mte, Mt & Marguerite \textup{(Calmels)} & 43\\
C\LJsup{tis}, Curt\LJsup{is}~; C\LJsup{zt} & Curtis~; Crouzet & 6\\
W, W., W\LJsup{m}, W\LJsup{–} & Williams \textup{(capitaine)} & 9\\
Pant & Panton \textup{(?), capitaine non identifié} & 4\\
Hippte~; Mad\LJsup{ne}~; Ls~; Cuis\LJsup{r} & Hippolyte~; Madeleine~; Louis \textup{(XI)}~; Cuisinier & 4\\
B\LJsup{te}~; Guilh & Baptiste \textup{(?)}~; Guilhem \textup{(?)} & 2\\
B~; Tl\LJsup{s}, Toul~; St Ser. & Bailleul~; Toulouse~; Saint-Sernin & 4\\
Ach, Achx~; P.\LJsup{v}~; Pas.~; Miraumt~; d'Hér\LJsup{t} & Acheux~; Puchevillers~; Pas-en-Artois~; Miraumont~; Hérissart & 6\\
Achiet-le-P\LJsup{t}.~; Molliens-au-B~; SFl & Achiet-le-Petit~; Molliens-au-Bois~; Saint-Flochel & 3\\
Luxemb~; Langdc~; V., V. Hugo~; Bournmth & Luxembourg~; Languedoc \textup{(Thermes du)}~; Victor-Hugo \textup{(marché)}~; Bournemouth & 5\\
\addlinespace[3pt]
\multicolumn{3}{@{}l}{\normalfont\bfseries Anglais abrégé, sigles et unités}\\*
\addlinespace[1pt]
Coy, Coy., Cy & \textit{Company} \textup{(compagnie britannique)} & 20\\
A.S.C., A.S.C & \textit{Army Service Corps} \textup{(train des équipages)} & 5\\
H.Q., HQ\LJsup{rs} & \textit{Headquarters} \textup{(quartier général)} & 3\\
R.T.O. & \textit{Railway Transport Officer} \textup{(commissaire de gare)} & 2\\
N.C.O., N.C.O's~; Q.M.S & \textit{Non-Commissioned Officer(s)}~; \textit{Quartermaster-Serjeant} & 4\\
M.P.s~; M.O~; M.T.~; G.S. & \textit{Military Police}~; \textit{Medical Officer}~; \textit{Mechanical} \textup{(ou} \textit{Motor}\textup{)} \textit{Transport}~; \textit{General Service} & 4\\
R.F.A., R.F.A~; Batt\LJsup{t} F.A–~; R. Fus\LJsup{rs}~; Bedfs & \textit{Royal Field Artillery}~; \textit{Battery, Field Artillery}~; \textit{Royal Fusiliers}~; \textit{Bedfords} & 5\\
Exp. F.~; Army P. Board & \textit{Expeditionary Force (Canteen)}~; \textit{Army Purchase Board} & 2\\
2\LJsup{nd}, 13th, 17\LJsup{th}, 18\LJsup{th}, 63\LJsup{rd}~; V\LJsup{th} & \textit{second, thirteenth…}~; \textit{Fifth} \textup{(V\LJsup{e} armée)} & 6\\
div, div\LJsup{on} \textnormal{(81\LJsup{st} div, 63\LJsup{rd} div)} & \textit{Division} & 3\\
Y.M.C.A., Y.M.C.A, Y.M-C-A- & \textit{Young Men's Christian Association} \textup{(foyer du soldat)} & 14\\
handk\LJsup{ch}~; p. \textnormal{(1 p. socks)} & \textit{handkerchief} \textup{(mouchoir)}~; \textit{pair} & 2\\
D.O.S.T. & \textup{sigle non attesté ; sans doute} \textit{Director} \textup{ou} \textit{Deputy Director of Supplies and Transport} \textup{(?)} & 1\\
miss. fr.~; M. F\LJsup{se} & Mission française & 2\\
artrie~; infrie~; cavrie & artillerie~; infanterie~; cavalerie & 4\\
R.A.T.~; G.V.~; P.L.M.~; W.C & réserve de l'armée territoriale~; grande vitesse~; Paris-Lyon-Méditerranée~; \textit{water-closet} & 4\\
A.P.C~; R2 & \textup{sigles non résolus} & 2\\
\end{longtable}
\endgroup
